\documentclass[10pt,a4paper]{amsart}
\usepackage[margin=19mm]{geometry}
\usepackage{amsmath,amssymb,mathtools}
\usepackage[scr=rsfs,cal=euler]{mathalfa}
\usepackage{xfrac}
\usepackage[unicode,hidelinks,bookmarks=false]{hyperref}
\newcommand{\ie}{i.e.\ }
\newcommand{\cf}{cf.\ }
\newcommand{\resp}{respectively\ }
\newcommand{\Subsection}{\subsection}
\providecommand{\nobreakdash}{\nobreak\hbox{-}}
\setlength{\emergencystretch}{2em}
\multlinegap0pt
\usepackage{bm}
\usepackage{bbm}
\usepackage{dsfont}
\usepackage{xspace}
\usepackage{cancel}
\usepackage{mathtools}
\usepackage{cleveref}
\usepackage{amsthm}
\makeatletter
\@ifundefined{mythm}{\newtheorem{mythm}{Mythm}}{}
\@ifundefined{mylemma}{\newtheorem{mylemma}[mythm]{Lemma}}{}
\makeatother
\DeclareMathOperator*{\colim}{colim}
\newcommand{\Cech}{\v{C}ech }
\newcommand{\RR}{\mathbb{R}}
\newcommand{\ZZip}{\ZZ[p^{-1}]}
\newcommand{\ZZ}{\mathbb{Z}}
\newcommand{\pRoots}{{1/p^\infty}}
\title[Perfect-Prismatic F-Crystals and p-adic Shtukas in Families]{Perfect-Prismatic F-Crystals and p-adic Shtukas in Families}
\author{Anton G\"uthge}
\date{Source: arXiv:2307.01108v3 (CC BY 4.0)}
\begin{document}
\maketitle

\noindent\emph{Source and licence.} Perfect-Prismatic F-Crystals and p-adic Shtukas in Families. Version arXiv:2307.01108v3 (CC BY 4.0), distributed under the Creative Commons Attribution 4.0 International licence, as recorded in the arXiv licence metadata for this exact version and shown on the abstract page https://arxiv.org/abs/2307.01108v3. Full rights notice: licenses/arxiv-2307.01108-CC-BY-4.0.txt.\par\smallskip

\noindent\textit{Editorial scope.} Counted unit: the Mylemma whose source label is \texttt{KimExtLaurentLemma} in the article (source0.tex lines 1190--1206, source label \texttt{KimExtLaurentLemma}), delivered together with the article's own complete original proof of it (source0.tex lines 1208--1274, one \texttt{proof} environment of 67 lines). No mathematics is written, repaired or re-derived: statement and proof are the article's own text line for line. Required context retained uncounted: KimExtCechCalculation (lines 1276--1284). Every definition, display and label of those lines is retained unchanged. Omitted from this package and not redistributed: 1--1189, 1207--1207, 1275--1275, 1285--2848. Nothing inside a retained excerpt is omitted or altered: every retained body is delivered exactly as the article's lines are written, so no omission receipt of any kind accompanies this package. Citations: the retained excerpts carry no citation command at all, so no bibliography entry of the article is transcribed and no reference is redistributed. Reference adaptation: none was needed and none was made; every reference inside the delivered unit resolves inside it, so no unresolved reference or citation is produced. Printed slips kept literal: no printed word, symbol, hypothesis or number of the counted statement or of its proof was altered. Layout and class: the article's own class is \texttt{article} with packages including fontenc, inputenc, babel, amsmath, amsthm, amssymb, newtxtext, newtxmath, mathtools, thmtools, tikz-cd, graphicx, adjustbox, imakeidx; the wrapper is the lane's pinned \texttt{amsart} editorial wrapper at 10pt on A4, which restates the theorem-like environments the retained text needs and adds the article's own macro definitions that the retained text needs (\texttt{\textbackslash Cech}, \texttt{\textbackslash RR}, \texttt{\textbackslash ZZ}, \texttt{\textbackslash ZZip}, \texttt{\textbackslash colim}, \texttt{\textbackslash pRoots}), with extra wrapper packages (bm, bbm, dsfont, xspace, cancel, mathtools, cleveref). The delivered numbering is the unit's own. Assets: no retained span contains a figure environment, a graphics-inclusion command, a PDF-page inclusion, a table or a TikZ picture; no image file is copied into this package, the package declares no asset dependency and no image file is redistributed with it. Licence: version arXiv:2307.01108v3 of this article is distributed under the Creative Commons Attribution 4.0 International licence, as recorded in the arXiv licence metadata for this exact version and shown on the abstract page https://arxiv.org/abs/2307.01108v3. A full rights notice is delivered at licenses/arxiv-2307.01108-CC-BY-4.0.txt. Characters: every retained excerpt is pure ASCII, so the delivered engine, XeLaTeX with its default Latin Modern text font, reports no missing character.
\begin{mylemma}\label{KimExtCechCalculation}
    The \Cech complex
    \[C_\infty\colon \quad0 \longrightarrow R^+((T_0^{\pRoots}))\longrightarrow  R^+((T_0^{\pRoots}, T_1^{\pRoots}))\longrightarrow\cdots\]
    with differentials defined by
    \begin{align*}\tag{**}\label{KimExtDifferential}
        T_0^{d_0}\cdots T_n^{d_n}\mapsto \sum_{i=0}^{n+1} (-1)^i T_0^{d_0}\cdots T_{i-1}^{d_{i-1}}\cdot T_{i+1}^{d_i}\cdots T_{n+1}^{d_n}
    \end{align*}
    is acyclic in positive degrees.
\end{mylemma}

\begin{mylemma}\label{KimExtLaurentLemma}
    Let \(P\) be the set of partial functions \(\ZZip\to \ZZip\), i.e., pairs \((D,m)\) where \(D\subset \ZZip\) and \(m\colon D\to \ZZip, \; d\mapsto m_d\). We will generally say \(D=D_m\) and regard it as part of the datum of \(m\). Let \(S\) be the subset
    \[\left\{m\in P
    \;\middle|\;\substack{m_d\leq \min(0, \frac dn ) \text{ for each \(d\in D_m\)},\\
    \#\left\{d\in D_m \;\mid\;  \sum_{i=1}^n s_i\left(m_d + \delta_{ij}(d-nm_d)\right)<C \text{ for some } j\in \{1,\dots,n\}\right\} \\
    \text{ is finite for all \(C\in \RR\), \(s_1,\dots,s_n\in \RR_{>0}.\)}
    }
    \right\}.\]
    We define a partial order on \(P\) (and hence on \(S\)) by 
    \[m\leq m' \quad \iff \quad D_m\subseteq D_{m'} \emph{ and } \forall d\in D_m\colon m_d\geq m_d'.\]
    Note the change of directions in the inequalities. Then there is an isomorphism of \(R^+\)-modules:
    \[R^+((T_1^{\pRoots},\dots, T_n^{\pRoots}))=\colim_{m\in S} \;\;\prod_{d\in D_m}
    \;\bigoplus_{\substack{d_1,\dots,d_n\in \ZZ[p^{-1}]\\
    \forall i : d_i\geq m_d\\
    d_1+\cdots+d_n=d}}
    R^+ T_1^{d_1}\cdots T_n^{d_n}\]
\end{mylemma}

\begin{proof}
    First, let's remark on the definition of \(S\): For any pair \((d,m_d)\) the set of tuples \((d_1,\dots,d_n)\) over which we take the direct sum forms a scaled and shifted \(n-1\)-standard-simplex lying in the (\(\sum d_i = d\))-plane. Indeed, at least if \(m_d\leq \frac dn\) (which we enforce in the definition of S), it is the convex hull of the scaled unit vectors 
    \[(d-nm_d)e_1,\dots, (d-nm_d)e_n\in \ZZip^n,\]
    shifted by \(m_d\sum_{i=1}^n e_i\). So through this lens, some \(m\in P\) can be seen as a set of disjoint shifted standard simplices in \(\ZZip^n\), indexed by \(d\in D_m\). The partial order is defined so that \(m\leq m'\) if and only if the set of shifted simplices defined by \(m\) is contained in the set of shifted simplices defined by \(m'\) (as subsets of \(\ZZip^n\)).\par
    Now if any point \((d_1,\dots,d_n)\in \ZZip^n\) of a shifted simplex defined by \((d, m_d)\) lies on the negative side of a hyperplane \(H\) specified by \(s_1,\dots,s_n,C\), then by convexity the same must be already true for one of its \(n\) corner points 
    \[\sum_i(m_d + \delta_{ij}(d-nm_d))e_i.\] 
    Thus for an \(m\in P\) to be in \(S\), for any hyperplane \(H\) defined by \(s_1,\dots,s_n,C\), only finitely many simplices may have a corner point (or, equivalently, any point) lying on the negative side of \(H\). The condition \(m_d\leq \frac 1n d\) guarantees that the simplices are non-empty, while \(m_d\leq 0\) enforces that the (shifted) simplex is large enough not to be fully contained in the \(d_i>0\)-quadrant, a technical assumption that will be useful in \Cref{KimExtCechCalculation} - this result would also be valid without it.\par
    Now the two descriptions of \(R^+((T^{\pRoots}_1,\dots,T_n^{\pRoots}))\) that we want to compare in this lemma clearly lie in the set of \emph{all} formal sums
    \[\sum_{\mathbf{d}=(d_1,\dots,d_n)\in \ZZip^n}a_{\mathbf{d}}^{}T_1^{d_1}\cdots T_n^{d_n},\]
    so we can show double inclusion. We will use the shorthand \(|\mathbf d|\coloneqq \sum_{i=1}^n d_i\).
    \begin{itemize}
        \item["\(\supseteq\)":]  Every element \(a\) in the group on the right side of the isomorphism lies in
        \[\prod_{d\in D_m}
        \;\bigoplus_{\substack{d_1,\dots,d_n\in \ZZ[p^{-1}]\\
        d_1,\dots,d_n\geq m_d\\
        |\mathbf d|=d}}
        R^+ T_1^{d_1}\cdots T_n^{d_n}\]
        for some \(m\in S\). Let \(H\) be an oriented hyperplane defined by \(s_1,\dots,s_n,C\). Then by the discussion about \(S\) above only finitely many \(d\in D_m\) define shifted simplices having at least one point lying on the negative side of \(H\), and any such \((d,m_d)\) can contribute only finitely many summands, so in total, \(a\) has only finitely many non-zero summands \(a_{\mathbf{d}}^{}T_1^{d_1}\cdots T_n^{d_n}\) such that \(\mathbf{d}\) lies on the negative side of \(H\). Thus \(a\in R^+((T_1^{\pRoots},\dots,T_n^{\pRoots}))\).
        \item["\(\subseteq\)":] Assume 
        \[a=\sum_{\mathbf{d}\in \ZZip^n} a_{\mathbf{d}}^{}T_1^{d_1}\cdots T_n^{d_n}\in R^+((T_1^{\pRoots},\dots,T_n^{\pRoots})).\]
        We define 
        \[D_m\coloneqq\{|\mathbf d|\;:\; a_{\mathbf d}\neq 0\}\]
        and for each \(d\in D_m\)
        \[m_d\coloneqq \min\Big(0,\;
        \min_{\substack{i\in\{1,\dots,n\}\\\mathbf d\in \ZZip^n\\|\mathbf d|=d,\;a_{\mathbf d}\neq 0}} d_i\Big).\]
        This automatically enforces \(m_d\leq 0\) and \(m_d\leq \frac dn\). The inner minimum is actually taken over a finite set: Indeed, by the definition of \(R^+((T_1^{\pRoots},\dots, T_n^{\pRoots}))\), there are only finitely many non-zero summands \(a_{\mathbf d}T_1^{d_1}\cdots T_n^{d_n}\) for which \(|\mathbf d|<d+1\) (by choosing \(s_1=\cdots=s_n=1, C= d+1\)). In particular, \(m_d\) is well-defined. This finiteness also shows that
        \[a\in\prod_{d\in D_m}
        \;\;\bigoplus_{\substack{d_1,\dots,d_n\in \ZZ[p^{-1}]\\
        d_i\geq m_d\\
        |\mathbf d|=d}}
        R^+ T_1^{d_1}\cdots T_n^{d_n}.\]
        Now all that is left to show is \(m\in S\)\footnote{The relevant picture here is the following: A priori it could happen that for some \(a\in R^+((T_1^{\pRoots},\dots, T_n^{\pRoots}))\) the subset \(\{\mathbf d \mid a_{\mathbf d}\neq 0\}\subset \ZZip^n\) has only finitely many points on the negative side of every oriented hyperplane defined as above, but that once we take the ``simplicial hull" (as in, the set of simplices defined by \(m\colon D_m\to \RR\)) of \(\{\mathbf d \mid a_{\mathbf d}\neq 0\}\), those simplices have infinitely many corners on the negative side of some hyperplane. We have to show that this can, in fact, not happen}. Assume \(m\notin S\), i.e., that there is some choice of a hyperplane defined by \(C\in \RR,s_1,\dots,s_n\in \RR_{>0}\) and an infinite subset \(D_<\subseteq D_m\) such that for each \(d\in D_<\) there is some \(j\in \{1,\dots,n\}\) for which
        \begin{align*}
            &\sum_{i=1}^n s_i\left(m_d + \delta_{ij}(d-nm_d)\right)\\
        =&m_d\left(\sum_{i=1}^n s_i \right)+ (d-nm_d)s_{j}
        <C\tag{*}\label{KimExtSupportBound}.
        \end{align*}
        First, as there are only finitely many choices for \(j\), there is some fixed \(j_0\) and some smaller but still infinite \(D_{<}'\subseteq D_{<}\) such that \eqref{KimExtSupportBound} holds true for \(j=j_0\) and each \(d\in D_{<}'\). Also, from the second form of the expression \eqref{KimExtSupportBound}, we see that the condition only depends on \(\sum_{i=1}^n s_i\) and \(s_{j_0}\), and that since \(d-nm_d>0\), decreasing \(s_{j_0}\) without altering \(\sum_{i=1}^n s_i\) will only make the expression smaller, so we can change the hyperplane so that \(s_{j_0}<\frac 1n\sum_{i=1}^n s_i\) and \eqref{KimExtSupportBound} still holds true.
        Rearranging the equation a different way
        \begin{align*}&\sum_{i=1}^n s_i\left(m_d + \delta_{ij_0}(d-nm_d)\right)\\=
        &s_{j_0}d+
        \left(-ns_{j_0}+\sum_{i=1}^ns_i\right)m_d<C
        \end{align*}
        we get
        \[m_d < -cd+b\]
        where
        \[c=s_{j_0}\left(-ns_{j_0}+\sum_{i=1}^ns_i\right)^{-1}, \qquad b=C\left(-ns_{j_0}+\sum_{i=1}^ns_i\right)^{-1},\]
        using that, by our assumption on \(s_{j_0}\) we have \(-ns_{j_0}+\sum_{i=1}^ns_i>0\) so we can divide by it without changing the direction of the inequality. The same assumption also shows \(c>0\). \par
        We have already used above that for any bound \(C'\) there are only finitely many \(\mathbf d\in \ZZip^n\) for which \(a_{\mathbf d} \neq 0\) and \(|\mathbf d|<C'\). Now for all \(d\geq C'\coloneqq\frac{b}{c}\), we have \(m_d < -cd+b \leq 0\), so only the (finitely many) \(d<C'\) can have \(m_d=0\). Thus, for each \(d\) in the infinite set \(D''_<\coloneqq D'_<\cap [C',\infty)\) we have 
        \[m_d=\min_{\substack{\mathbf d\in \ZZip^n\\|\mathbf d|=d,\;a_{\mathbf d}\neq 0\\i\in\{1,\dots,n\}}} d_i.\]
        Let \(\mathbf D''_< \subset \{\mathbf{d}\in\ZZip^n\mid a_\mathbf{d} \neq 0\}\) be the set of \(\mathbf d=(d_1,\dots, d_n)\) where that minimum is attained, i.e. 
        \[d=|\mathbf d|\in D''_<, \quad m_d = \min_{i\in \{1,\dots, n\}} d_i.\]
        Because there are only finitely many choices for \(i\) in the minimum above, we can assume that there is some \(i_0\) and some infinite subset \(\mathbf D'''_<\subseteq \mathbf D''_<\) such that for each \(\mathbf d\in \mathbf D'''_<\), the minimum is attained at \(i_0\), i.e. 
        \[d_{i_0}=m_d < -cd+b.\]
        Setting
        \[s'_i=\begin{cases}1 \quad &\text{ if } i\neq i_0\\1+\frac 1c&\text{ if } i=i_0\end{cases}, \qquad C'=\frac bc\]
        we get 
        \begin{align*}
            \sum_{i=1}^n s'_id_i
            &=\left(\sum_{i=1}^n d_i \right)+ \frac{d_{i_0}}c\\
            &< d+\frac{(-cd+b)}c\\
            &=\frac b c \\
            &=C'
        \end{align*}
        for each of the infinitely many \(\mathbf d\in \mathbf{D}'''_<\), each of which has \(a_{\mathbf d}\neq 0\). This shows that \(a\notin R^+((T_1^{\pRoots},\dots, T_n^{\pRoots}))\).
    \end{itemize}
\end{proof}

\end{document}
